\documentclass[../DoAn.tex]{subfiles}
\begin{document}

Chương này sẽ giới thiệu chi tiết về các nền tảng lý thuyết và công nghệ cốt lõi được ứng dụng để xây dựng hệ thống giám sát người cao tuổi. Nhằm đảm bảo tính rõ ràng, nội dung được chia thành hai nhóm chính: nhóm công nghệ phần cứng và thuật toán nhận diện (phía thiết bị biên), và nhóm công nghệ phần mềm (phía máy chủ và ứng dụng quản lý).

\section{Cơ sở lý thuyết về Nhận diện hành vi và Phát hiện té ngã dựa trên cảm biến quán tính (IMU)}
\label{section:3.1}

\subsection{Cơ chế nhận diện hành vi con người (HAR)}
Nhận diện hành vi con người (Human Activity Recognition - HAR) thông qua cảm biến quán tính (IMU) hoạt động dựa trên nguyên lý thu thập và phân tích các mẫu động học của cơ thể. Mỗi hành động như đi bộ (Walk), chạy (Run), hay đứng yên (Idle) đều tạo ra một mẫu tín hiệu (pattern) gia tốc và vận tốc góc đặc trưng, mang tính tuần hoàn và lặp lại trên không gian 3 chiều.
\begin{itemize}
    \item \textbf{Đứng yên (Idle):} Trục gia tốc dọc theo chiều trọng lực của Trái Đất (thường là trục Z hoặc Y tùy vị trí đeo) sẽ đo được giá trị xấp xỉ $1g$ ($9.8 m/s^2$), trong khi các trục còn lại xấp xỉ $0g$. Vận tốc góc từ con quay hồi chuyển gần như bằng 0.
    \item \textbf{Đi bộ (Walk) và Chạy (Run):} Bước chân tạo ra các xung gia tốc có tính chu kỳ hình sin. Đi bộ có biên độ dao động và tần số thấp, trong khi chạy tạo ra các đỉnh gia tốc (peaks) sắc nhọn với biên độ lớn và tần số cao hơn. Con quay hồi chuyển cũng ghi nhận sự vung tay hoặc chuyển động xoay của cơ thể tương ứng với từng bước đi.
\end{itemize}

Bằng cách cắt tín hiệu liên tục thành các \textit{cửa sổ trượt (sliding windows)} với kích thước cố định (ví dụ: 2 giây tương đương 200 mẫu ở tần số 100Hz), chúng ta thu được các ma trận dữ liệu hai chiều chứa toàn bộ thông tin không gian và thời gian của hành động. Các mô hình học máy và học sâu sẽ tự động trích xuất các đặc trưng tiềm ẩn từ ma trận này để phân loại chính xác hành vi \cite{nazar2025csohm, booranawong2025stairs}.

\subsection{Cơ chế phát hiện té ngã (Fall Detection)}
Khác với các hành vi mang tính chu kỳ, té ngã là một chuỗi sự kiện bất thường, diễn ra cực kỳ nhanh và có đặc điểm động lực học chia làm 4 giai đoạn rõ rệt:
\begin{enumerate}
    \item \textbf{Giai đoạn trước khi ngã (Pre-fall):} Người bệnh đang thực hiện các hoạt động bình thường.
    \item \textbf{Giai đoạn rơi tự do (Free-fall):} Cơ thể bắt đầu mất thăng bằng và rơi xuống. Trong khoảnh khắc này, gia tốc tổng hợp (Sum Vector Magnitude - SVM) của cả 3 trục tiến dần về mức $0g$ do trạng thái không trọng lượng tạm thời.
    \item \textbf{Giai đoạn va chạm (Impact):} Đây là đặc trưng quan trọng nhất. Cơ thể đập vào mặt đất sinh ra một lực phản hồi cực lớn. Gia tốc kế sẽ ghi nhận một đỉnh (peak) đột biến vọt lên rất cao, thường vượt qua mức $2.5g$ đến $3g$. Cùng lúc đó, con quay hồi chuyển ghi nhận một sự thay đổi góc nghiêng cực gắt.
    \item \textbf{Giai đoạn sau khi ngã (Post-fall):} Sau cú va chạm, cơ thể nằm bất động trên mặt sàn. Dữ liệu gia tốc kế sẽ trở lại trạng thái tĩnh, nhưng lúc này sự phân bố trọng trường trên các trục không gian đã thay đổi hoàn toàn so với tư thế đứng (do cơ thể nằm ngang).
\end{enumerate}

Nhờ sự kết hợp của cả \textbf{đỉnh gia tốc đột biến} và \textbf{sự thay đổi tư thế tĩnh (orientation change)}, hệ thống có thể dễ dàng phân biệt được một cú té ngã nguy hiểm so với các hành động mạnh khác như nhảy lên, hoặc việc bệnh nhân chủ động nằm xuống giường \cite{tseng2025wearable, lee2021deep}.

\section{Nhóm Công nghệ Phần cứng và Giao thức mạng}
\label{section:3.2}

\subsection{Vi điều khiển ESP32-S3 (Module Seeed Studio XIAO)}
ESP32-S3 là một hệ thống trên chip (SoC) mạnh mẽ do Espressif Systems phát triển \cite{espressif2023esp32s3}, được tích hợp vi xử lý lõi kép Xtensa 32-bit LX7 hoạt động ở xung nhịp lên đến 240 MHz. Điểm nổi bật nhất của ESP32-S3 so với các phiên bản tiền nhiệm là việc được bổ sung bộ tập lệnh mở rộng SIMD (Single Instruction, Multiple Data). Nhờ đó, chip có khả năng tính toán song song các mảng dữ liệu lớn, giúp gia tăng đáng kể tốc độ xử lý mạng nơ-ron thông qua thư viện ESP-NN \cite{espressif2023espnn}. 

Trong đồ án này, để tối ưu hóa kích thước cho thiết bị đeo (wearable), hệ thống sử dụng module \textbf{Seeed Studio XIAO ESP32S3} \cite{seeed2023xiao}. Đây là một bo mạch phát triển có kích thước siêu nhỏ (chỉ cỡ ngón tay cái) nhưng vẫn mang đầy đủ sức mạnh của chip ESP32-S3, với 512 KB SRAM tích hợp và 8MB bộ nhớ PSRAM (phù hợp cho các tác vụ cần bộ nhớ đệm lớn mà cần tốc độ xử lý cao như stream video MJEPG, tensor arena cho tinyML \cite{warden2019tinyml}). Vi điều khiển này đóng vai trò là "bộ não" trung tâm, không chỉ làm nhiệm vụ thu thập tín hiệu từ các cảm biến ngoại vi qua giao thức I2C, mà còn trực tiếp thực thi mô hình học máy (TinyML) ngay tại thiết bị biên (Edge Computing). 

Một ưu điểm vượt trội khác của module XIAO ESP32S3 là thiết kế phần cứng thân thiện với thiết bị đeo di động. Ở mặt sau của bo mạch được tích hợp sẵn 2 pad hàn (chân BAT+ và BAT-) dành riêng cho việc kết nối trực tiếp với pin Lithium 3.7V. Tính năng này kết hợp với mạch quản lý sạc pin đã được tích hợp sẵn trên bo mạch giúp thiết bị có thể hoạt động hoàn toàn độc lập, dễ dàng được cung cấp nguồn điện trực tiếp từ pin mà không cần phải thiết kế thêm mạch nguồn cồng kềnh.

\begin{figure}[H]
    \centering
    \includegraphics[width=1\textwidth]{Hinhve/esp32s3_front.png}
    \includegraphics[width=1\textwidth]{Hinhve/esp32s3_back.png}
    \caption{Mặt trước và mặt sau của module Seeed Studio XIAO ESP32S3 (chú ý 2 chân BAT cấp nguồn pin ở mặt sau)}
    \label{fig:esp32s3}
\end{figure}

\subsection{Cảm biến quán tính MPU6050}
MPU6050 là một cảm biến vi cơ điện tử (MEMS) 6 bậc tự do (6-DOF) bao gồm một gia tốc kế (Accelerometer) 3 trục và một con quay hồi chuyển (Gyroscope) 3 trục \cite{invensense2013mpu6050}. 
\begin{itemize}
    \item \textbf{Gia tốc kế:} Đo lường lực gia tốc tuyến tính (bao gồm cả gia tốc trọng trường), giúp xác định sự thay đổi vận tốc và hướng của thiết bị. Trong bài toán phát hiện té ngã, gia tốc kế chịu trách nhiệm phát hiện các xung lực va chạm cực đại khi người bệnh tiếp đất.
    \item \textbf{Con quay hồi chuyển:} Đo lường vận tốc góc của thiết bị theo 3 trục không gian. Thông số này cực kỳ hữu ích trong việc tính toán góc nghiêng của cơ thể, hỗ trợ phân biệt các tư thế sinh hoạt bình thường và trạng thái mất thăng bằng.
\end{itemize}
Dữ liệu thô từ MPU6050 được lấy mẫu liên tục thông qua giao tiếp I2C với tốc độ cao, tạo thành các chuỗi dữ liệu đầu vào cho mạng nơ-ron \cite{Paramasivam2024Heliyon}.

\begin{figure}[H]
    \centering
    \includegraphics[width=0.45\textwidth]{Hinhve/mpu6050.jpg}
    \caption{Cảm biến quán tính MPU6050 đo gia tốc và vận tốc góc}
    \label{fig:mpu6050}
\end{figure}

\subsection{Module truyền thông 4G LTE A7680C}
Nhằm khắc phục những hạn chế của kết nối Wi-Fi (phụ thuộc vào cơ sở hạ tầng mạng cục bộ) và Bluetooth (yêu cầu phải ghép nối với điện thoại thông minh), đồ án sử dụng module SIM 4G LTE A7680C \cite{simcom2021a7680c}. Đây là một module viễn thông siêu nhỏ gọn, hỗ trợ băng tần LTE Cat-1, tối ưu hóa cho các ứng dụng IoT (Internet of Things) đòi hỏi độ trễ thấp và băng thông vừa phải. Nhờ tích hợp module A7680C, thiết bị đeo có thể hoạt động hoàn toàn độc lập, liên tục truyền tải dữ liệu cảnh báo và viễn trắc trực tiếp lên máy chủ đám mây từ bất kỳ vị trí nào có phủ sóng di động \cite{tseng2025wearable}.

\begin{figure}[H]
    \centering
    \includegraphics[width=0.45\textwidth]{Hinhve/a7680c_front.jpeg}
    \includegraphics[width=0.45\textwidth]{Hinhve/a7680c_back.jpeg}
    \caption{Module truyền thông mạng di động 4G LTE A7680C (Mặt trước và mặt sau)}
    \label{fig:a7680c}
\end{figure}

\subsection{Giao thức truyền thông PPPoS}
Để kết nối module 4G A7680C với vi điều khiển ESP32-S3 một cách tối ưu, đồ án không sử dụng tập lệnh AT truyền thống để gửi dữ liệu từng bước, mà ứng dụng giao thức PPPoS (Point-to-Point Protocol over Serial) \cite{lwip2023pppos}. PPPoS cho phép thiết lập một kết nối mạng trực tiếp qua cổng nối tiếp (UART). Sau khi kết nối thành công, hệ điều hành FreeRTOS trên ESP32 sẽ nhận diện kết nối 4G như một giao diện mạng (Network Interface) thực thụ, có địa chỉ IP riêng. Việc này giúp các ứng dụng tầng trên (như MQTT client hay HTTP client) có thể tương tác với mạng di động bằng các socket tiêu chuẩn một cách mượt mà và ổn định.

\begin{figure}[H]
    \centering
    \includegraphics[width=0.8\textwidth]{Hinhve/pppos_stack.png}
    \caption{Kiến trúc giao thức PPPoS ánh xạ mạng 4G thành giao diện mạng trên ESP32S3}
    \label{fig:pppos_stack}
\end{figure}

\subsection{Giao thức truyền tải MQTT}
MQTT (Message Queuing Telemetry Transport) là một giao thức nhắn tin theo mô hình Xuất bản/Đăng ký (Publish/Subscribe), được thiết kế đặc biệt cho các thiết bị IoT hoạt động trong điều kiện băng thông thấp và mạng không ổn định \cite{oasis2019mqtt}. Khác với kiến trúc Request/Response nặng nề của HTTP, MQTT sử dụng một máy chủ trung gian (Broker) để điều phối thông điệp. 
\begin{itemize}
    \item \textbf{Cơ chế QoS (Quality of Service):} MQTT hỗ trợ 3 mức QoS (0, 1, 2). Trong đồ án, các thông điệp cảnh báo té ngã khẩn cấp được đặt mức QoS 1 (At least once) để đảm bảo gói tin chắc chắn được máy chủ tiếp nhận, trong khi dữ liệu viễn trắc thông thường dùng QoS 0 để tối ưu băng thông.
    \item \textbf{Keep-Alive:} MQTT duy trì một kết nối TCP liên tục nhưng tiêu tốn cực ít năng lượng thông qua cơ chế gửi gói tin Ping định kỳ, giúp máy chủ lập tức nhận biết khi thiết bị mất tín hiệu hoặc cạn pin.
\end{itemize}

\begin{figure}[H]
    \centering
    \includegraphics[width=1\textwidth]{Hinhve/mqtt_architechture.png}
    \caption{Mô hình giao tiếp Publish/Subscribe của MQTT trong hệ thống giám sát}
    \label{fig:mqtt_architecture}
\end{figure}

\section{Công nghệ Nhận diện hành vi (TinyML)}
\label{section:3.3}

\subsection{Mạng nơ-ron tích chập thặng dư 1 chiều (1D ResNet)}
Trong lĩnh vực xử lý chuỗi thời gian (time-series) như tín hiệu IMU, mạng nơ-ron tích chập 1 chiều (1D CNN) tỏ ra vượt trội so với các mạng nơ-ron hồi quy (RNN/LSTM) về tốc độ suy luận (inference speed) và khả năng trích xuất đặc trưng cục bộ. Tuy nhiên, khi mạng CNN được thiết kế quá sâu nhằm tăng độ chính xác, hiện tượng suy biến đạo hàm (vanishing gradient) sẽ xuất hiện. 

Để giải quyết vấn đề này, kiến trúc 1D ResNet (Residual Network) được ứng dụng \cite{he2016resnet}. Điểm cốt lõi của ResNet là cơ chế \textit{Skip Connection} (kết nối tắt), cho phép dữ liệu từ các tầng trước "bỏ qua" một số tầng trung gian để cộng trực tiếp vào đầu ra của các tầng sau. Cơ chế này không chỉ giúp huấn luyện các mạng cực sâu một cách hiệu quả mà còn giữ lại được các thông tin dao động tĩnh quan trọng của cảm biến, vốn rất dễ bị mất mát sau nhiều phép tích chập.

\subsection{Kỹ thuật lượng tử hóa (Quantization) cho TinyML}
TinyML (Trí tuệ nhân tạo vạn vật) là khái niệm chỉ việc triển khai các mô hình học máy lên các vi điều khiển có tài nguyên phần cứng cực kỳ eo hẹp (vài chục KB RAM) \cite{warden2019tinyml}. Một mô hình 1D ResNet được huấn luyện bằng TensorFlow/Keras thường biểu diễn trọng số dưới dạng số thực dấu phẩy động 32-bit (Float32). Dạng dữ liệu này có kích thước rất lớn và yêu cầu vi xử lý tính toán số thực mạnh mẽ, gây tốn pin và vượt quá giới hạn SRAM của ESP32.

Để giải quyết rào cản này, kỹ thuật lượng tử hóa sau huấn luyện (Post-Training Quantization - PTQ) được áp dụng. Toàn bộ trọng số và tham số kích hoạt của mạng nơ-ron được ánh xạ từ không gian Float32 xuống số nguyên 8-bit (INT8). Nhờ đó, kích thước mô hình (ROM) giảm đi 4 lần (từ vài MB xuống chỉ còn dưới 100 KB), đồng thời vi điều khiển ESP32-S3 có thể tận dụng tập lệnh SIMD để thực hiện phép nhân ma trận số nguyên siêu tốc. Lượng tử hóa tuy có làm giảm một phần nhỏ độ chính xác của mô hình, nhưng hoàn toàn nằm trong mức độ chấp nhận được để đổi lấy khả năng chạy trực tiếp trên thiết bị biên.

\section{Nền tảng Phần mềm và Hệ thống}
\label{section:3.4}

\subsection{Nền tảng phát triển Nhúng ESP-IDF và FreeRTOS}
ESP-IDF (Espressif IoT Development Framework) là bộ công cụ phát triển phần mềm chính thức cho dòng chip ESP32 \cite{espressif2023espidf}. Khác với môi trường Arduino IDE vốn chỉ phù hợp cho các dự án nhỏ với cấu trúc lập trình tuần tự (Super Loop), ESP-IDF cung cấp khả năng can thiệp sâu vào phần cứng và hỗ trợ đa luồng thông qua hệ điều hành thời gian thực FreeRTOS \cite{freertos2023manual}.

FreeRTOS đóng vai trò là "nhạc trưởng" điều phối toàn bộ hoạt động của thiết bị biên. Thay vì bắt bộ vi xử lý thực hiện các công việc nối tiếp nhau và phải chờ đợi (blocking), FreeRTOS cho phép chia nhỏ chương trình thành các tác vụ (Tasks) độc lập. Trong đồ án này, các tác vụ như lấy mẫu cảm biến IMU, nội suy AI, và truyền thông MQTT được cấp phát các mức độ ưu tiên (priority) khác nhau. Hệ thống sử dụng cơ chế hàng đợi (Queue) và cờ hiệu (Semaphore) để đồng bộ hóa dữ liệu giữa các tác vụ, đảm bảo rằng cảm biến không bao giờ bị bỏ lỡ mẫu dữ liệu nào ngay cả khi mô hình AI đang chiếm dụng CPU để tính toán.

\subsection{Nền tảng triển khai mô hình học máy (TinyML Pipeline)}
Đồ án xây dựng một quy trình (pipeline) khép kín để phát triển và triển khai mô hình học máy lên thiết bị nhúng với các công cụ chuyên dụng:
\begin{itemize}
    \item \textbf{Huấn luyện mô hình (Training):} Quá trình xây dựng và huấn luyện mạng 1D ResNet được thực hiện trên máy tính bằng thư viện mã nguồn mở \textbf{TensorFlow và Keras (phiên bản 2.21)} \cite{abadi2016tensorflow}. Đây là framework mạnh mẽ giúp định nghĩa các lớp tích chập, hàm mất mát và tối ưu hóa trọng số mô hình với độ chính xác cao.
    \item \textbf{Nén và Lượng tử hóa (Quantization):} Sau khi huấn luyện, mô hình được chuyển đổi thông qua \textbf{TensorFlow Lite (TFLite)}. Công cụ này thực hiện lượng tử hóa sau huấn luyện (Post-Training Quantization), ép kiểu trọng số từ Float32 xuống INT8, giúp giảm dung lượng mô hình tới 4 lần mà gần như không làm suy giảm độ chính xác.
    \item \textbf{Suy luận tại biên (Inference):} Mô hình đã lượng tử hóa được nạp vào ESP32-S3 và thực thi bằng thư viện \textbf{TensorFlow Lite for Microcontrollers} \cite{david2021tensorflowlite}. Đây là một phiên bản rút gọn của TFLite, được thiết kế đặc biệt để chạy trên vi điều khiển mà không cần hệ điều hành nền tảng máy chủ, tiêu tốn chỉ vài chục kilobyte bộ nhớ.
    \item \textbf{Tối ưu hóa phần cứng (Hardware Acceleration):} Để tối đa hóa tốc độ suy luận, đồ án tích hợp thư viện \textbf{ESP-NN} \cite{espressif2023espnn} do Espressif cung cấp. Thư viện này chứa các hàm toán học (như Convolution, Fully Connected) được viết bằng hợp ngữ (Assembly) và tối ưu hóa chuyên biệt cho tập lệnh SIMD của lõi Xtensa trên ESP32-S3, giúp quá trình nhận diện diễn ra gần như tức thời.
\end{itemize}

\subsection{Bộ lọc Kalman 1D cho Dung hợp Cảm biến (Sensor Fusion)}

Trong bài toán ước lượng góc nghiêng (Roll, Pitch) của thiết bị đeo từ tín hiệu IMU, mỗi loại cảm biến đều mang một điểm mạnh và điểm yếu đối nghịch nhau. Gia tốc kế cho phép tính toán góc tĩnh rất chính xác dựa trên hướng của véc-tơ trọng lực, nhưng lại cực kỳ nhạy cảm với nhiễu tần số cao (high-frequency noise) khi thiết bị va đập mạnh hoặc chuyển động đột ngột --- đúng là đặc trưng của cú ngã. Ngược lại, con quay hồi chuyển (gyroscope) phản ứng tức thời và mượt mà với mọi chuyển động, nhưng khi ước lượng góc bằng cách tích phân vận tốc góc theo thời gian, sai số nhỏ của mỗi chu kỳ lấy mẫu sẽ tích lũy thành sai số trôi (drift) không thể tự hiệu chỉnh được.

Bộ lọc Kalman 1 chiều (1D Kalman Filter) giải quyết mâu thuẫn này bằng cách kết hợp linh hoạt đầu ra của cả hai cảm biến theo một trọng số tối ưu được cập nhật động theo từng chu kỳ \cite{kalman1960new}. Thuật toán hoạt động qua hai bước tuần hoàn: bước \textit{Dự báo (Predict)} sử dụng vận tốc góc từ gyroscope để ước lượng góc nghiêng tại thời điểm hiện tại và tính toán sai số dự báo tích lũy; bước \textit{Cập nhật (Update)} sử dụng góc tĩnh tính từ gia tốc kế như một phép đo tham chiếu để hiệu chỉnh lại kết quả dự báo thông qua hệ số Kalman Gain. Tham số phương sai nhiễu đo lường của gia tốc kế (\texttt{R\_measure}) đóng vai trò điều chỉnh mức độ tin tưởng vào phép đo gia tốc: giá trị cao hơn sẽ làm bộ lọc ít bị dao động khi gia tốc nhiễu (như lúc va chạm), trong khi giá trị thấp hơn giúp hội tụ nhanh hơn khi cơ thể ở trạng thái tĩnh.

Trong hệ thống này, thư viện \texttt{lib\_kalman} được gọi trực tiếp trong \texttt{svc\_imu} tại mỗi chu kỳ lấy mẫu 100~Hz. Góc Roll và Pitch đã được làm mượt bởi bộ lọc được đưa vào cùng với dữ liệu gia tốc thô và vận tốc góc thô để tạo thành véc-tơ đặc trưng 6 chiều đầu vào cho mô hình TinyML. Điều này giúp loại bỏ phần lớn nhiễu ngắn hạn không mang thông tin hành vi, đồng thời cung cấp cho mô hình một biểu diễn tư thế ổn định hơn, từ đó giảm thiểu đáng kể tỷ lệ cảnh báo sai trong sinh hoạt hàng ngày.

\section{Hệ quản trị Cơ sở dữ liệu}
\label{section:3.5}

Hệ thống ứng dụng kiến trúc cơ sở dữ liệu kép (Dual-Database) để xử lý hiệu quả hai loại dữ liệu có tính chất hoàn toàn khác biệt.

\subsection{Cơ sở dữ liệu chuỗi thời gian InfluxDB}
Dữ liệu thu thập từ các cảm biến IoT (gia tốc, số bước chân, mức pin) mang tính chất chuỗi thời gian (time-series), tức là được sinh ra liên tục, số lượng cực lớn và chỉ được ghi nối tiếp chứ hiếm khi bị sửa đổi. Đối với loại dữ liệu này, các cơ sở dữ liệu quan hệ truyền thống sẽ nhanh chóng bị quá tải. 

Do đó, đồ án sử dụng \textbf{InfluxDB}, một hệ quản trị cơ sở dữ liệu chuỗi thời gian chuyên biệt. InfluxDB lưu trữ dữ liệu theo cấu trúc Measurement (bảng), Tags (thẻ phân loại có đánh chỉ mục, ví dụ: \texttt{device\_id}), và Fields (giá trị thực tế, ví dụ: \texttt{accX, accY}). Với cơ chế nén dữ liệu cực tốt và ngôn ngữ truy vấn Flux mạnh mẽ, InfluxDB có thể đáp ứng hàng chục nghìn luồng ghi (write) mỗi giây và truy xuất dữ liệu lịch sử cực nhanh để vẽ biểu đồ giám sát.

\subsection{Cơ sở dữ liệu quan hệ PostgreSQL}
Ngược lại với dữ liệu cảm biến, các thông tin về hồ sơ bệnh nhân, danh sách thiết bị, phân quyền y tá và lịch sử cảnh báo lại có tính quan hệ chặt chẽ và yêu cầu tính toàn vẹn dữ liệu (ACID) tuyệt đối. Để giải quyết bài toán này, đồ án sử dụng \textbf{PostgreSQL} - hệ quản trị cơ sở dữ liệu quan hệ mã nguồn mở mạnh mẽ nhất hiện nay. PostgreSQL xử lý việc liên kết (JOIN) các bảng dữ liệu phức tạp, đảm bảo rằng một thiết bị tại một thời điểm chỉ được gán cho một bệnh nhân duy nhất (Unique Constraint), và duy trì tính nhất quán cho toàn bộ hệ thống quản trị. Đặc biệt, PostgreSQL được thiết kế theo kiến trúc \textbf{Đa khách thuê (Multi-tenancy)}, bổ sung trường \texttt{org\_id} vào các bảng nhằm cách ly dữ liệu an toàn giữa các viện dưỡng lão khác nhau.

% \begin{figure}[H]
%     \centering
%     % \includegraphics[width=0.8\textwidth]{Hinhve/dual_database.png}
%     \caption{Mô hình Cơ sở dữ liệu kép kết hợp InfluxDB và PostgreSQL}
%     \label{fig:dual_database}
% \end{figure}

\section{Nền tảng Web và Dịch vụ Backend}
\label{section:3.6}

\subsection{Framework phát triển Backend FastAPI}
Trái tim của máy chủ hệ thống được xây dựng bằng \textbf{FastAPI}, một web framework hiện đại và hiệu năng cao dành cho Python \cite{ramirez2023fastapi}.
\begin{itemize}
    \item \textbf{Hiệu năng cao và Bất đồng bộ (Async):} FastAPI hỗ trợ xử lý bất đồng bộ (Asynchronous) gốc thông qua ASGI. Điều này cực kỳ quan trọng khi Backend phải đồng thời lắng nghe hàng trăm thông điệp từ MQTT Broker, ghi vào InfluxDB và phục vụ các HTTP Request từ Frontend mà không bị chặn luồng (blocking).
    \item \textbf{Xác thực dữ liệu tự động:} Nhờ tích hợp chặt chẽ với thư viện Pydantic, FastAPI tự động kiểm tra tính hợp lệ của các gói tin JSON gửi lên, giảm thiểu đáng kể lỗi do dữ liệu không đúng định dạng.
    \item \textbf{Tài liệu API tự sinh:} FastAPI tự động tạo ra tài liệu API tương tác trực quan (Swagger UI), giúp quá trình kiểm thử và tích hợp với Frontend diễn ra trơn tru, nhanh chóng.
\end{itemize}

% \begin{figure}[H]
%     \centering
%     % \includegraphics[width=0.8\textwidth]{Hinhve/fastapi_swagger.png}
%     \caption{Giao diện tài liệu API Swagger tự sinh từ FastAPI}
%     \label{fig:fastapi_swagger}
% \end{figure}

\subsection{Framework Frontend Next.js và Tailwind CSS}
Giao diện quản lý dành cho nhân viên y tế (Dashboard) được xây dựng bằng \textbf{Next.js} (phiên bản App Router) kết hợp cùng công cụ đóng gói \textbf{Turbopack}, một framework phát triển dựa trên thư viện ReactJS.
\begin{itemize}
    \item \textbf{Next.js:} Khác với ReactJS thuần túy (chỉ render ở phía client), Next.js cung cấp khả năng kết hợp giữa Server-Side Rendering (SSR) và Client-Side Rendering (CSR) \cite{vercel2023nextjs}. Nhờ đó, ứng dụng tải trang cực nhanh, tối ưu trải nghiệm người dùng và giúp bảo mật các khóa truy cập API quan trọng ở phía máy chủ.
    \item \textbf{Quản lý trạng thái (State Management):} Hệ thống ứng dụng \textbf{TanStack Query} để tối ưu hóa việc gọi API, quản lý bộ nhớ đệm (caching) và đồng bộ dữ liệu với máy chủ, kết hợp với \textbf{Zustand} để quản lý các trạng thái giao diện cục bộ một cách nhẹ nhàng.
    \item \textbf{Tailwind CSS và Shadcn UI:} Đây là một framework CSS theo phong cách \textit{Utility-first} kết hợp cùng thư viện UI components hiện đại \cite{tailwind2023css}. Thay vì phải viết các file CSS độc lập, lập trình viên áp dụng trực tiếp các class định dạng ngay trong JSX. Công cụ này giúp tạo ra một giao diện Web trực quan, độ tùy biến cao, hỗ trợ thiết kế đáp ứng (Responsive Design) mượt mà.
\end{itemize}

% \begin{figure}[H]
%     \centering
%     % \includegraphics[width=0.8\textwidth]{Hinhve/nextjs_architecture.png}
%     \caption{Cơ chế Server-Side Rendering (SSR) trong framework Next.js}
%     \label{fig:nextjs_architecture}
% \end{figure}

\end{document}